\documentclass[10pt,openany]{book}
% One editorial system for the five existing Studium collections.
% Load after \documentclass, before book-specific configuration.
\usepackage[utf8]{inputenc}
\usepackage[T1]{fontenc}
\usepackage{lmodern}
\DeclareUnicodeCharacter{2212}{\ensuremath{-}}
\DeclareUnicodeCharacter{2192}{\ensuremath{\to}}
\DeclareUnicodeCharacter{03B1}{\ensuremath{\alpha}}
\usepackage[spanish,es-nodecimaldot]{babel}
\spanishplainpercent
\AtBeginDocument{\spanishdeactivate{"~<>}}
\usepackage[a4paper,margin=2.5cm,headheight=14pt]{geometry}
\usepackage{amsmath,amssymb,mathtools,booktabs,array,longtable,graphicx,xcolor,tikz,enumitem,microtype,titlesec,fancyhdr,listings,etoolbox}
\usetikzlibrary{arrows.meta,calc,positioning,patterns}
\usepackage[breakable,skins]{tcolorbox}
\usepackage[labelfont=bf]{caption}
\usepackage[hyphens]{url}
\usepackage[hidelinks,unicode,hypertexnames=false]{hyperref}
\hypersetup{pdfauthor={Daniel Miguel Tejedor}}
\addto\captionsspanish{\renewcommand{\contentsname}{Índice general}\renewcommand{\tablename}{Cuadro}}
\AtBeginDocument{\normalfont\DeclareFontShape{T1}{lmr}{bx}{sc}{<->ssub*lmr/m/sc}{}}
\definecolor{studiumInk}{HTML}{142C40}
\definecolor{studiumRule}{HTML}{142C40}
\definecolor{studiumBlue}{RGB}{24,74,130}
\definecolor{studiumCyan}{RGB}{0,184,205}
\definecolor{studiumLight}{RGB}{232,240,250}
\definecolor{studiumTipBack}{RGB}{232,244,236}
\definecolor{studiumTipFrame}{RGB}{27,107,58}
\definecolor{studiumDefBack}{RGB}{232,240,250}
\definecolor{studiumDefFrame}{RGB}{24,74,130}
\definecolor{studiumWorkedBack}{RGB}{255,244,230}
\definecolor{studiumWorkedFrame}{RGB}{138,78,12}
\definecolor{studiumCheckBack}{RGB}{243,236,248}
\definecolor{studiumCheckFrame}{RGB}{96,48,130}
\definecolor{studiumCodeBack}{RGB}{245,245,242}
\definecolor{studiumCodeRule}{RGB}{48,48,52}
\titleformat{\chapter}[display]{\normalfont\filright}{\large\mdseries\scshape\chaptertitlename\ \thechapter}{1ex}{\normalfont\huge\bfseries}
\titlespacing*{\chapter}{0pt}{2.5ex plus 1ex minus .2ex}{2.3ex}
\titleformat{name=\chapter,numberless}[display]{\normalfont\filright}{}{0pt}{\normalfont\huge\bfseries}
\titleformat{\section}{\normalfont\Large\bfseries}{\thesection}{.8em}{}
\titleformat{\subsection}{\normalfont\large\bfseries}{\thesubsection}{.7em}{}
\titleformat{\subsubsection}{\normalfont\normalsize\bfseries}{\thesubsubsection}{.7em}{}
\makeatletter
\renewcommand*\l@section{\@dottedtocline{1}{1.5em}{3.5em}}
\@openrightfalse
\let\cleardoublepage\clearpage
\makeatother
\setcounter{secnumdepth}{2}\setcounter{tocdepth}{1}
\setlength{\parindent}{1.5em}\setlength{\parskip}{3pt plus .5pt}
\setlength{\emergencystretch}{2em}\raggedbottom
\setlist{itemsep=2pt,topsep=4pt,leftmargin=*}
\setlist[description]{style=nextline,leftmargin=0pt,labelwidth=0pt,labelindent=0pt}
\allowdisplaybreaks[2]
\newcounter{studiumfigure}
\makeatletter\@addtoreset{studiumfigure}{chapter}\makeatother
\renewcommand{\thestudiumfigure}{\thechapter.\arabic{studiumfigure}}
\newcommand{\studiumcaption}[2]{\refstepcounter{studiumfigure}\label{#1}\addcontentsline{lof}{figure}{\protect\numberline{\thestudiumfigure}#2}{\small\textbf{Figura~\thestudiumfigure.} #2\par}}
\newcommand{\Raire}{287.05287}
\newcommand{\gzero}{9.80665}
\tcbset{studiumcommon/.style={breakable,boxrule=.65pt,arc=.7pt,left=1.7mm,right=1.7mm,top=1.3mm,bottom=1.3mm,before skip=10pt,after skip=10pt,fonttitle=\bfseries,coltitle=white,coltext=black},
 studiumtip/.style={studiumcommon,colback=studiumTipBack,colframe=studiumTipFrame,colbacktitle=studiumTipFrame},
 studiumdef/.style={studiumcommon,colback=studiumDefBack,colframe=studiumDefFrame,colbacktitle=studiumDefFrame},
 studiumworked/.style={studiumcommon,colback=studiumWorkedBack,colframe=studiumWorkedFrame,colbacktitle=studiumWorkedFrame},
 studiumcheck/.style={studiumcommon,colback=studiumCheckBack,colframe=studiumCheckFrame,colbacktitle=studiumCheckFrame}}
\newtcolorbox{studiumworked}[1]{studiumworked,title={#1}}
\newcommand{\studiumchapterguide}[1]{\par\noindent\textcolor{studiumInk}{\textbf{Orientación del capítulo.}}\enspace\textit{#1}\par\medskip}
\newcommand{\objetivos}[1]{\studiumchapterguide{#1}}
\newcommand{\caja}[2]{\begin{tcolorbox}[studiumdef,title={#1}]#2\end{tcolorbox}}
\newcommand{\resumen}[1]{\begin{tcolorbox}[studiumcheck,title={Síntesis del capítulo}]#1\end{tcolorbox}}
\newcommand{\studiumchaptercheck}{\begin{tcolorbox}[studiumcheck,title={Comprobación final}]
\begin{enumerate}
\item Explica el modelo del capítulo y las hipótesis que permiten utilizarlo.
\item Reproduce un desarrollo o un ejemplo sin consultar su solución y comprueba unidades, signos y condiciones de validez.
\item Resuelve las preguntas o ejercicios disponibles y distingue el resultado calculado de su interpretación física.
\end{enumerate}\end{tcolorbox}}
\newcommand{\figura}[2]{\begin{figure}[tbp]\centering\includegraphics[width=.83\linewidth]{#1.pdf}\caption{#2}\label{fig:#1}\end{figure}\noindent La figura~\ref{fig:#1} ofrece una representación del modelo discutido.}
\newsavebox{\studiumMathBox}
\newcommand{\fitmath}[1]{\sbox{\studiumMathBox}{$\displaystyle #1$}\ifdim\wd\studiumMathBox>\linewidth\resizebox{\linewidth}{!}{\usebox{\studiumMathBox}}\else\usebox{\studiumMathBox}\fi}
\pagestyle{fancy}\fancyhf{}
\fancyhead[LE]{\footnotesize\thepage\hspace{1.5em}\mdseries\scshape\nouppercase{\leftmark}}
\fancyhead[RO]{\footnotesize\mdseries\scshape\nouppercase{\leftmark}\hspace{1.5em}\thepage}
\renewcommand{\headrulewidth}{.4pt}\renewcommand{\footrulewidth}{0pt}
\fancypagestyle{plain}{\fancyhf{}\fancyfoot[C]{\small\thepage}\renewcommand{\headrulewidth}{0pt}}
\lstdefinestyle{studium}{basicstyle=\ttfamily\small,columns=fullflexible,keepspaces=true,showstringspaces=false,breaklines=true,breakatwhitespace=false,frame=leftline,framerule=1.15pt,rulecolor=\color{studiumCodeRule},backgroundcolor=\color{studiumCodeBack},xleftmargin=1.15em,framexleftmargin=.75em,aboveskip=.9em,belowskip=.7em,tabsize=4,keywordstyle=\ttfamily\bfseries,commentstyle=\ttfamily\itshape,stringstyle=\ttfamily}
\lstset{style=studium}
% Shared cover: every letter remains editable vector text; photographs keep original pixels.
\providecommand{\studiumcoverphoto}{}
\providecommand{\studiumcoveredition}{Primera edición}
\providecommand{\studiumcovercategory}{MANUAL UNIVERSITARIO}
\providecommand{\studiumphotocredit}{Crédito y procedencia en el proyecto del libro.}
\providecommand{\studiumphototop}{0}
\newsavebox{\studiumPhotoBox}
\newcommand{\studiumcover}[3]{%
\hypersetup{pageanchor=false}\begin{titlepage}\thispagestyle{empty}
\begin{tikzpicture}[remember picture,overlay,x=1pt,y=-1pt]
\begin{scope}[shift={(current page.north west)}]
\fill[studiumInk] (0,0) rectangle (595.276,841.89);
\begin{scope}
\clip (0,89.5) rectangle (595.276,585.5);
\sbox{\studiumPhotoBox}{\includegraphics[width=\paperwidth]{\studiumcoverphoto}}%
\ifdim\ht\studiumPhotoBox<496pt
 \node[anchor=north,inner sep=0pt] at (297.638,89.5){\includegraphics[height=496pt]{\studiumcoverphoto}};
\else
 \node[anchor=north,inner sep=0pt] at (297.638,\studiumphototop){\usebox{\studiumPhotoBox}};
\fi
\end{scope}
\fill[white] (0,0) rectangle (595.276,89.5);
\node[anchor=west,text=studiumInk,font=\sffamily\bfseries\fontsize{18}{20}\selectfont,inner sep=0pt] at (51.223,38.5){STUDIUM};
\fill[studiumCyan] (51.223,61.5) rectangle (80.8,63.7);
\node[anchor=west,text=studiumInk,font=\sffamily\fontsize{9}{11}\selectfont,inner sep=0pt] at (89.49,65.5){BIBLIOTECA DE INGENIERÍA};
\node[anchor=west,text=studiumCyan,font=\sffamily\bfseries\fontsize{10}{12}\selectfont,inner sep=0pt] at (51.223,608){#3};
\node[anchor=north west,text=white,text width=493pt,font=\rmfamily\bfseries\fontsize{32}{35}\selectfont,inner sep=0pt] at (51.223,632){#1};
\node[anchor=north west,text=white!82!studiumInk,text width=493pt,font=\sffamily\fontsize{11}{14}\selectfont,inner sep=0pt] at (51.223,722){#2};
\node[anchor=west,text=white,font=\sffamily\bfseries\fontsize{10}{12}\selectfont,inner sep=0pt] at (51.223,763){GRADO EN INGENIERÍA AEROESPACIAL};
\node[anchor=west,text=white!82!studiumInk,font=\sffamily\fontsize{10}{12}\selectfont,inner sep=0pt] at (51.223,778){Universidad de León};
\draw[white!75!studiumInk,line width=.4pt] (51.223,792) -- (544.05,792);
\node[anchor=west,text=white!82!studiumInk,font=\sffamily\fontsize{10}{12}\selectfont,inner sep=0pt] at (51.223,815){2026};
\node[anchor=east,text=white,font=\sffamily\bfseries\fontsize{10}{12}\selectfont,inner sep=0pt] at (544.05,815){\studiumcoveredition};
\end{scope}\end{tikzpicture}\null\end{titlepage}\hypersetup{pageanchor=true}}
\newcommand{\portada}[2]{\studiumcover{#1}{#2}{\studiumcovercategory}}

\providecommand{\studiumvolumetype}{manual}
\newcommand{\studiumreadingguide}{%
\chapter*{Cómo utilizar este volumen}\addcontentsline{toc}{chapter}{Cómo utilizar este volumen}
\ifdefstring{\studiumvolumetype}{manual}{\noindent\textbf{Ruta de estudio.} Empieza por la orientación del capítulo y sus conceptos. Escribe las hipótesis antes de utilizar las ecuaciones y sigue los desarrollos con sus unidades. Después intenta los problemas y las preguntas de repaso sin consultar las soluciones. Cierra cada capítulo con la comprobación final.}{%
\ifdefstring{\studiumvolumetype}{problemas}{\noindent\textbf{Ruta de trabajo.} Lee el enunciado y copia sus datos con unidades. Define las incógnitas, el modelo y las hipótesis. Intenta resolver el problema antes de consultar el desarrollo. Compara después el resultado, la interpretación física y las comprobaciones.}{\noindent\textbf{Ruta de entrenamiento.} Lee las instrucciones particulares de cada simulacro y reserva el tiempo propuesto. Trabaja primero con los enunciados. Después utiliza las soluciones indicadas en el volumen para revisar el planteamiento, los cálculos y las comprobaciones.}}

\noindent\textbf{Recuadros.} El azul identifica definiciones y resultados; el verde, consejos; el naranja, problemas resueltos; y el violeta, síntesis y autoevaluación. Los apartados técnicos conservan los nombres propios de cada asignatura.

\noindent\textbf{Problemas y simulacros.} Separa enunciado, datos, hipótesis, modelo, desarrollo, resultado y comprobación. En los simulacros trabaja primero con los enunciados y consulta después las soluciones. Sus tiempos y puntuaciones son propuestas de entrenamiento cuando así se indican, sin sustituir las instrucciones oficiales del curso.

\noindent\textbf{Consulta.} El índice general permite localizar los temas. La notación, los formularios, las soluciones y la bibliografía mantienen sus apartados al final del manual cuando forman parte de él.
\section*{Créditos y alcance editorial}\addcontentsline{toc}{section}{Créditos y alcance editorial}
\noindent Daniel Miguel Tejedor. Colección Studium. Edición de estudio independiente para el Grado en Ingeniería Aeroespacial, Universidad de León.

\noindent La compilación y la unificación editorial no añaden una aprobación institucional ni una nueva verificación científica. Cada libro conserva las hipótesis, límites y comprobaciones de sus modelos; los informes registran su alcance y sus pendientes.

\noindent\textbf{Imagen de portada.} \studiumphotocredit
}

% Book-specific metadata; all visual rules live in shared/latex/.
\renewcommand{\studiumcoverphoto}{../figures/cover.jpeg}
\renewcommand{\studiumcoveredition}{Primera edición}
\renewcommand{\studiumphotocredit}{Fotografía conservada del PDF de clase aportado por el usuario. La identificación y el crédito de origen siguen pendientes.}
\renewcommand{\studiumphototop}{0}
\graphicspath{{../figures/}}

\hypersetup{pdftitle={Turbomáquinas}}
\renewcommand{\studiumvolumetype}{simulacros}
\begin{document}
\studiumcover{TURBOMÁQUINAS}{Entrenamiento y criterios de resolución}{SIMULACROS}
\frontmatter
\studiumreadingguide
\tableofcontents
\mainmatter
\chapter*{Presentación y alcance}
Propuesta original de estudio, sin carácter oficial. Cada simulacro dispone de 120 minutos orientativos. Declare convenios, hipótesis y unidades; conserve los signos tangenciales. Las soluciones están separadas en la colección de problemas, con los mismos identificadores.\par\medskip\noindent Puntuación orientativa: A1 y A2, 4 puntos cada uno; A3, 2 puntos. B1, 4 puntos; B2 y B3, 3 puntos cada uno. En cada ejercicio: planteamiento y convenio, 25 por ciento; desarrollo, 35; sustitución y unidades, 25; interpretación y comprobaciones, 15.\par\chapter{Simulacro A}
\section{C19-01-P}
\(C19-01-P \cdot A1 \cdot\) Cierre de compresor axial. Un rotor compresor al radio medio recibe Cx y remolinos dados. Obtenga trabajo, temperatura total final, presión total final y área geométrica de entrada. Resultado final: área de entrada.
\begin{center}\begin{tabular}{p{.15\linewidth}p{.42\linewidth}p{.31\linewidth}}\textbf{Símbolo} & \textbf{Significado} & \textbf{Valor}\\\hline
\(T_{0e}\) & temperatura total entrada & \(302\,\mathrm{K}\)\\
\(p_{0e}\) & presión total entrada & \(118000\,\mathrm{Pa}\)\\
\(C_x\) & componente axial & \(156\,\mathrm{m}\,\mathrm{s}^{-1}\)\\
\(C_{\theta e}\) & remolino entrada rotor & \(38\,\mathrm{m}\,\mathrm{s}^{-1}\)\\
\(C_{\theta s}\) & remolino salida rotor & \(147\,\mathrm{m}\,\mathrm{s}^{-1}\)\\
\(U\) & velocidad periférica & \(291\,\mathrm{m}\,\mathrm{s}^{-1}\)\\
\(\dot m\) & caudal & \(21\,\mathrm{kg}\,\mathrm{s}^{-1}\)\\
\(\eta_{\mathit{tt}}\) & eficiencia total-total & \(0.87\)\\
\(B\) & fracción disponible & \(0.96\)\\
\(c_p\) & calor específico & \(1004.5\,\mathrm{J}\,\mathrm{kg}^{-1}\,\mathrm{K}^{-1}\)\\
\(\gamma\) & razón de calores & \(1.4\)\\
\(R_g\) & constante gas & \(287\,\mathrm{J}\,\mathrm{kg}^{-1}\,\mathrm{K}^{-1}\)\\
\end{tabular}\end{center}
\noindent\textbf{Hipótesis del caso:} Gas perfecto, radio medio y Cx constantes, escalón adiabático.\par
\noindent\textbf{Comprobación exigida:} indique el dominio físico, reconstruya al menos un balance y explique qué datos faltarían para evaluar una máquina real.\par
\begin{center}\begin{tikzpicture}\draw[thick] (0,0)--(7,0);\draw[thick] (0,2.5)--(7,2.5);\node[draw,fill=orange!20,minimum width=2cm,minimum height=1.5cm] at (3.5,1.25) {Rotor};\draw[->,blue,thick] (.2,1.2)--(2,1.2) node[midway,above] {$C_x,C_{\theta e}$};\draw[->,blue,thick] (5,1.2)--(6.8,1.2) node[midway,above] {$C_x,C_{\theta s}$};\draw[->,orange,thick] (3.5,2.8)--(3.5,3.6) node[above] {$U>0$};\end{tikzpicture}\end{center}
\clearpage
\section{C19-02-P}
\(C19-02-P \cdot A2 \cdot\) Turbina con descarga axial. Diseñe un escalón repetitivo de turbina que descarga axialmente. Calcule carga, trabajo, temperatura total de salida y presión total de salida.
\begin{center}\begin{tabular}{p{.15\linewidth}p{.42\linewidth}p{.31\linewidth}}\textbf{Símbolo} & \textbf{Significado} & \textbf{Valor}\\\hline
\(T_{0e}\) & temperatura total entrada & \(1490\,\mathrm{K}\)\\
\(p_{0e}\) & presión total entrada & \(940000\,\mathrm{Pa}\)\\
\(\mathcal R\) & reacción & \(0.48\)\\
\(U\) & velocidad periférica & \(314\,\mathrm{m}\,\mathrm{s}^{-1}\)\\
\(\eta_{\mathit{tt}}\) & eficiencia total-total & \(0.895\)\\
\(c_p\) & calor específico & \(1156.696969697\,\mathrm{J}\,\mathrm{kg}^{-1}\,\mathrm{K}^{-1}\)\\
\(\gamma\) & razón de calores & \(1.33\)\\
\end{tabular}\end{center}
\noindent\textbf{Hipótesis del caso:} Repetitivo, mismo radio y Cx constantes, sin refrigeración.\par
\noindent\textbf{Comprobación exigida:} indique el dominio físico, reconstruya al menos un balance y explique qué datos faltarían para evaluar una máquina real.\par
\begin{center}\begin{tikzpicture}\node[draw,fill=blue!15,minimum height=1.5cm] (s) at (2,1) {Estator};\node[draw,fill=orange!20,minimum height=1.5cm] (r) at (5,1) {Rotor};\draw[->] (0,1)--(s);\draw[->] (s)--(r) node[midway,above] {2};\draw[->] (r)--(7,1) node[right] {$C_{\theta3}=0$};\node at (0,1.6) {1};\node at (7,.4) {3};\draw[->,orange,thick] (5,2)--(5,2.8) node[above] {$U$};\end{tikzpicture}\end{center}
\clearpage
\section{C19-03-P}
\(C19-03-P \cdot A3 \cdot\) Distribución radial y tracción. Un escalón compresor tiene leyes libres en ambas estaciones. Calcule trabajo radialmente uniforme, reacción en raíz y tensión centrífuga nominal de pala uniforme. Resultado final: tensión.
\begin{center}\begin{tabular}{p{.15\linewidth}p{.42\linewidth}p{.31\linewidth}}\textbf{Símbolo} & \textbf{Significado} & \textbf{Valor}\\\hline
\(\Omega\) & velocidad angular & \(1090\,\mathrm{s}^{-1}\)\\
\(K_e\) & constante entrada & \(14\,\mathrm{m}^{2}\,\mathrm{s}^{-1}\)\\
\(K_s\) & constante salida & \(44\,\mathrm{m}^{2}\,\mathrm{s}^{-1}\)\\
\(r_h\) & radio raíz & \(0.218\,\mathrm{m}\)\\
\(r_t\) & radio punta & \(0.302\,\mathrm{m}\)\\
\(\rho\) & densidad material & \(8050\,\mathrm{kg}\,\mathrm{m}^{-3}\)\\
\end{tabular}\end{center}
\noindent\textbf{Hipótesis del caso:} Reacción reducida de escalón repetitivo; pala mecánica uniforme.\par
\noindent\textbf{Comprobación exigida:} indique el dominio físico, reconstruya al menos un balance y explique qué datos faltarían para evaluar una máquina real.\par
\begin{center}\begin{tikzpicture}\draw[fill=gray!15,thick] (1,0) rectangle (2,3.7);\draw[thick] (.5,0)--(2.5,0);\foreach \x in {.6,.9,...,2.4}{\draw (\x,0)--(\x-.25,-.25);}\draw[->,orange] (1.5,2)--(1.5,4.5) node[above] {Inercia centrífuga};\draw[<->] (3.2,0)--(3.2,3.7) node[midway,right] {$r_t-r_h$};\node[left] at (1,0) {$r_h$};\node[left] at (1,3.7) {$r_t$};\node at (5.5,2.8) {$rC_\theta=K_v$};\node at (5.5,1.7) {$U=\Omega r$};\end{tikzpicture}\end{center}
\chapter{Simulacro B}
\section{C19-04-P}
\(C19-04-P \cdot B1 \cdot\) Rodete radial y prestación. Un rodete de salida radial sin prerrotación tiene radio y giro dados. Con slip y eficiencia prescritos, halle trabajo, temperatura total y relación de presión.
\begin{center}\begin{tabular}{p{.15\linewidth}p{.42\linewidth}p{.31\linewidth}}\textbf{Símbolo} & \textbf{Significado} & \textbf{Valor}\\\hline
\(\Omega\) & velocidad angular & \(2070\,\mathrm{s}^{-1}\)\\
\(r\) & radio salida & \(0.206\,\mathrm{m}\)\\
\(\sigma_s\) & factor Ctheta/Ctheta ideal & \(0.875\)\\
\(\eta_{\mathit{tt}}\) & eficiencia total-total & \(0.825\)\\
\(T_{0e}\) & temperatura total entrada & \(296\,\mathrm{K}\)\\
\(c_p\) & calor específico & \(1004.5\,\mathrm{J}\,\mathrm{kg}^{-1}\,\mathrm{K}^{-1}\)\\
\(\gamma\) & razón de calores & \(1.4\)\\
\end{tabular}\end{center}
\noindent\textbf{Hipótesis del caso:} Gas perfecto, salida radial ideal antes de slip; sin prerrotación.\par
\noindent\textbf{Comprobación exigida:} indique el dominio físico, reconstruya al menos un balance y explique qué datos faltarían para evaluar una máquina real.\par
\begin{center}\begin{tikzpicture}\draw[thick] (2.5,1.7) circle (1.5);\draw[thick] (2.5,1.7) circle (.45);\foreach \a in {0,45,...,315}{\draw[blue,rotate around={\a:(2.5,1.7)}] (2.95,1.7)--(4,1.7);}\draw[->,orange,thick] (3,3.6) arc (70:140:1.8);\node at (2.5,4) {$\Omega>0$};\draw[->] (4,1.7)--(6,1.7) node[right] {$C_{r2}$};\draw[->] (4,1.7)--(4,2.8) node[right] {$C_{\theta2}$};\node at (5.7,.5) {Prerrotación nula};\end{tikzpicture}\end{center}
\clearpage
\section{C19-05-P}
\(C19-05-P \cdot B2 \cdot\) Cascada y elección entera. Una fila compresora tiene velocidades relativas y cambio tangencial dados. Se ha elegido un número entero de álabes; calcule paso, solidez real y difusión.
\begin{center}\begin{tabular}{p{.15\linewidth}p{.42\linewidth}p{.31\linewidth}}\textbf{Símbolo} & \textbf{Significado} & \textbf{Valor}\\\hline
\(r\) & radio & \(0.279\,\mathrm{m}\)\\
\(c\) & cuerda real & \(0.054\,\mathrm{m}\)\\
\(Z_b\) & número de álabes & \(43\)\\
\(V_e\) & módulo relativo entrada & \(296\,\mathrm{m}\,\mathrm{s}^{-1}\)\\
\(V_s\) & módulo relativo salida & \(224\,\mathrm{m}\,\mathrm{s}^{-1}\)\\
\(|\Delta V_\theta|\) & magnitud de cambio tangencial & \(103\,\mathrm{m}\,\mathrm{s}^{-1}\)\\
\end{tabular}\end{center}
\noindent\textbf{Hipótesis del caso:} Indicador de difusión definido en referencia relativa.\par
\noindent\textbf{Comprobación exigida:} indique el dominio físico, reconstruya al menos un balance y explique qué datos faltarían para evaluar una máquina real.\par
\begin{center}\begin{tikzpicture}\foreach \y in {0,1.8,3.6}{\draw[blue,thick] (0,\y) .. controls (1,\y+.8) and (2.8,\y+.8) .. (4,\y+.3);\draw[blue,thick] (0,\y) .. controls (1,\y+.2) and (2.8,\y+.3) .. (4,\y+.3);}\draw[<->] (-.5,0)--(-.5,1.8) node[midway,left] {$s$};\draw[<->] (0,4.6)--(4,4.9) node[midway,above] {$c$};\node[align=center] at (6,2) {$Z_b$ entero\\$V_e,V_s,|\Delta V_\theta|$};\end{tikzpicture}\end{center}
\clearpage
\section{C19-06-P}
\(C19-06-P \cdot B3 \cdot\) Balance térmico integrado. Una pared se refrigera por convección interior. Calcule flujo por área, temperatura caliente y caudal mínimo energético para el incremento permitido. Resultado final: caudal.
\begin{center}\begin{tabular}{p{.15\linewidth}p{.42\linewidth}p{.31\linewidth}}\textbf{Símbolo} & \textbf{Significado} & \textbf{Valor}\\\hline
\(T_{aw}\) & temperatura recuperación exterior & \(1580\,\mathrm{K}\)\\
\(T_c\) & temperatura refrigerante local media del modelo & \(720\,\mathrm{K}\)\\
\(h_g\) & coeficiente exterior & \(1540\,\mathrm{W}\,\mathrm{m}^{-2}\,\mathrm{K}^{-1}\)\\
\(h_c\) & coeficiente interior & \(2320\,\mathrm{W}\,\mathrm{m}^{-2}\,\mathrm{K}^{-1}\)\\
\(t\) & espesor metal & \(0.0012\,\mathrm{m}\)\\
\(k_m\) & conductividad & \(23\,\mathrm{W}\,\mathrm{m}^{-1}\,\mathrm{K}^{-1}\)\\
\(A\) & superficie intercambio & \(0.028\,\mathrm{m}^{2}\)\\
\(c_p\) & calor específico refrigerante & \(1080\,\mathrm{J}\,\mathrm{kg}^{-1}\,\mathrm{K}^{-1}\)\\
\(\Delta T_c\) & incremento permitido & \(165\,\mathrm{K}\)\\
\end{tabular}\end{center}
\noindent\textbf{Hipótesis del caso:} Red estacionaria unidimensional, propiedades constantes, sin radiación.\par
\noindent\textbf{Comprobación exigida:} indique el dominio físico, reconstruya al menos un balance y explique qué datos faltarían para evaluar una máquina real.\par
\begin{center}\begin{tikzpicture}\node (a) at (0,1) {$T_{aw}$};\node[draw,minimum height=1cm] (b) at (1.8,1) {$1/h_g$};\node[draw,minimum height=1cm] (c) at (3.8,1) {$t/k_m$};\node[draw,minimum height=1cm] (d) at (5.8,1) {$1/h_c$};\node (e) at (7.3,1) {$T_c$};\draw[->] (a)--(b);\draw[->] (b)--(c);\draw[->] (c)--(d);\draw[->] (d)--(e);\node at (2.8,1.8) {$T_{w,g}$};\draw[->,red,thick] (1,2.5)--(6.5,2.5) node[midway,above] {$q^{\prime\prime}$};\end{tikzpicture}\end{center}
\end{document}
